% Source: https://artofproblemsolving.com/community/c3h2290582_finding_f_given_that_all_of_the_roots_of_a_polynomial_are_positive_reals
% Stable identity: aops-topic-2290582-q01
Let the roots, with multiplicity, be $r_1,\dots,r_8>0$. By Vieta's formulas,
\[\sum r_i=4,\qquad\sum_{i<j}r_ir_j=7,\qquad f=r_1r_2\cdots r_8\]
(the degree is even, so the constant term is $+r_1\cdots r_8$). Hence
\[\sum r_i^2=\left(\sum r_i\right)^2-2\sum_{i<j}r_ir_j=16-14=2.\]
By the Cauchy-Schwarz inequality, $\left(\sum r_i\right)^2\le8\sum r_i^2$, that is, $16\le16$. So equality holds, which happens only when all $r_i$ are equal: $r_i=\dfrac12$. Therefore $f=\left(\dfrac12\right)^8=\dfrac1{256}$.

This value does occur: $\left(x-\dfrac12\right)^8$ has $x^7$-coefficient $-8\cdot\dfrac12=-4$ and $x^6$-coefficient $\dbinom82\cdot\dfrac14=7$. So the only possible value is $\boxed{f=\dfrac{1}{256}}$.
